\section*{Chapitre \ref{ch:s1:index-rebalancing} --- Rééquilibrage des indices}

\begin{solution}{exo:s1:index-rebalancing:1}
Un non-membre entre au rang 285 ou mieux ($300 - 15$)~; un membre sort au-delà du rang 315 ($300 + 15$).
\end{solution}

\begin{solution}{exo:s1:index-rebalancing:2}
Les fonds doivent acheter $0.15/0.006 = 25$ jours de volume~: $0.7 \times 0.02 \times \sqrt{25} = 7.0~\%$.
\end{solution}

\begin{solution}{exo:s1:index-rebalancing:3}
La transaction d'annonce rapporte la moitié non prise en avance, $0.5 \times 7.0~\% = 3.5~\%$~; le fournisseur de liquidité gagne la moitié qui s'inverse, également 3.5~\%, avant
hors coûts et hors mouvements propres de l'action.
\end{solution}

\begin{solution}{exo:s1:index-rebalancing:4}
Une action proche du seuil de coupure doit franchir la bande pour changer de statut, de sorte que les petits mouvements autour du seuil ne la font pas entrer et sortir sans cesse~: moins de changements, donc moins de transactions
de la part des \omterm{def:s1:index-rebalancing:fund}{fonds indiciels}. Les changements deviennent plus rares et concernent des actions qui ont davantage bougé, ce qui les rend plus faciles à prévoir au \omterm{def:s1:index-rebalancing:fund}{jour de classement}, mais une bande fait aussi dépendre
le résultat du statut courant de l'action.
\end{solution}

\begin{solution}{exo:s1:index-rebalancing:5}
La demande des \omterm{def:s1:index-rebalancing:fund}{fonds indiciels}, et donc la poussée totale, est la même~; le capital d'arbitrage ne change que lorsque le cours bouge. Plus le capital prévoit le changement et
achète tôt, plus une part de la poussée se produit avant l'annonce et moins il en reste entre l'annonce et la clôture effective.
\end{solution}

\begin{solution}{exo:s1:index-rebalancing:6}
Environ un quart de ses positions portent sur des actions qui n'entrent pas et ne profitent d'aucune poussée (et peuvent même pâtir du débouclage des positions des autres)~; il manque aussi
plus de la moitié des ajouts effectifs. Son rendement par position est nettement inférieur à celui du prédicteur parfait, et son risque plus élevé.
\end{solution}

\begin{solution}{exo:s1:index-rebalancing:7}
Avec 30~\% des actions suivies et seulement un cinquième pris en avance, la transaction d'annonce rapporte 5.7~\% par événement~; avec 5~\% suivies et quatre cinquièmes pris en avance, 0.4~\%.
La première a une poussée importante (7.5~\% en moyenne) dont l'essentiel reste à venir~; la seconde une poussée faible (3.1~\%) dont l'essentiel a déjà disparu avant l'annonce.
\end{solution}

\begin{solution}{exo:s1:index-rebalancing:8}
L'événement a changé~: davantage de capital prévoit désormais les changements, donc une plus grande part du mouvement se produit avant l'annonce, et la part postérieure s'est réduite~;
Patel et Welch observent un retour à la moyenne dans les années 2000. Il faut mesurer l'événement sur les années récentes, avec les coûts et les dérives propres des actions, et se demander qui d'autre achète.
\end{solution}

\begin{solution}{pb:s1:index-rebalancing:1}
\begin{enumerate}
\item Un fonds qui réplique un indice dans ses pondérations~; la date à laquelle un indice fondé sur des règles mesure les critères d'appartenance~; la probabilité, avant le \omterm{def:s1:index-rebalancing:fund}{jour de classement},
qu'une action soit membre après la recomposition.
\item Semestriel à partir de 2026~: juin (\omterm{def:s1:index-rebalancing:fund}{jour de classement} le 30 avril, effet après la clôture du 26 juin) et décembre (\omterm{def:s1:index-rebalancing:fund}{jour de classement} le 30 octobre, effet après la clôture du
11 décembre).
\item Les 300 plus grandes actions, recomposées tous les 126 jours avec une bande de 15 rangs~; annonce 16 jours et clôture effective 40 jours après le \omterm{def:s1:index-rebalancing:fund}{jour de classement}~;
331 ajouts et 270 retraits sur dix-sept recompositions.
\item Il évite que les actions proches du seuil de coupure n'entrent et sortent sans cesse.
\item En projetant la capitalisation de chaque action avec sa propre volatilité, puis en appliquant la règle à chaque tirage.
\item 100~\% au \omterm{def:s1:index-rebalancing:fund}{jour de classement}~; 86~\% et 72~\% pour les ajouts 20 jours avant (retraits 86~\% et 74~\%)~; 76~\% et 45~\% soixante jours avant (76~\% et 46~\%).
\item Les capitalisations ont davantage de temps pour évoluer.
\item Des prévisions précoces laissent davantage de mouvement à venir, mais se trompent plus souvent.
\item Une part de la poussée avant l'annonce, le reste jusqu'à la clôture effective, dont la moitié s'inverse dans les 20 jours suivants.
\item 3.75~\%, 2.40~\% et 2.65~\% par événement, avec 15~\% des actions suivies et une moitié prise en avance.
\item Un rendement anormal significatif à l'annonce des ajouts au S\&P 500, lié aux achats des \omterm{def:s1:index-rebalancing:fund}{fonds indiciels}~; une hausse permanente pour les ajouts, mais aucune
baisse permanente pour les retraits.
\item Les actions qui entrent dans un indice sont des gagnantes récentes, et le marché synthétique crée du momentum.
\item \textbf{Résultat clé.} Les entrées et sorties d'indice sont prévues avec une précision de 86~\% et un rappel de 72~\% vingt jours avant le \omterm{def:s1:index-rebalancing:fund}{jour de classement} (76~\% et 45~\%
soixante jours avant)~; la transaction d'annonce rapporte de 2.2~\% à 5.7~\% par événement lorsque la part suivie passe de 5~\% à 30~\% avec un cinquième pris en avance, et de 0.4~\% à
1.3~\% avec quatre cinquièmes pris en avance.
\item Retour à la moyenne dans les années 2000~: pas de déplacement permanent de la demande pour les changements du S\&P 500.
\item Plus tôt, pour les prédicteurs, et plus tard, pour les fournisseurs de liquidité à la clôture effective.
\item Petit capital~: fourniture de liquidité à la clôture effective sur quelques valeurs~; grand capital~: prévisions sur l'ensemble de la liste.
\item Les déséquilibres des enchères de clôture, la demande attendue des \omterm{def:s1:index-rebalancing:fund}{fonds indiciels} et les cours des ajouts dont les positions sont les plus encombrées.
\item Deux événements plus petits par an au lieu d'un~: moins de demande par événement, plus d'événements, et un calendrier plus chargé pour les prédicteurs.
\item Sa date, sa taille et son sens sont connus à l'avance.
\item Une poussée connue et temporaire, dont l'essentiel est échangé avant et après la date par ceux qui la prévoient et ceux qui la subissent.
\end{enumerate}
\end{solution}

\begin{solution}{iq:s1:index-rebalancing:1}
Les \omterm{def:s1:index-rebalancing:fund}{fonds indiciels} doivent l'acheter immédiatement, et la contrepartie de cette demande n'est pas parfaitement élastique~: le cours monte jusqu'à ce que les vendeurs soient suffisamment rémunérés pour fournir les
actions. Une partie de la hausse peut relever de l'information ou de la prise de conscience~; le reste est une pression qui tend à s'inverser.
\end{solution}

\begin{solution}{iq:s1:index-rebalancing:2}
Prévoir les changements à partir des règles et des capitalisations avant le \omterm{def:s1:index-rebalancing:fund}{jour de classement}, acheter les probables ajouts et vendre les probables retraits avant la liste préliminaire, et
fournir de la liquidité à la clôture effective face à la demande des \omterm{def:s1:index-rebalancing:fund}{fonds indiciels}~; dimensionner la position selon la demande attendue rapportée au volume et selon la confiance dans la
prévision.
\end{solution}

\begin{solution}{iq:s1:index-rebalancing:3}
La capitalisation et le flottant à la date du jour, la volatilité de chaque action, les règles de l'indice, y compris les bandes et les filtres, et le calendrier~; simuler les capitalisations jusqu'au
\omterm{def:s1:index-rebalancing:fund}{jour de classement} pour obtenir une probabilité par action. L'incertitude croît avec le temps restant jusqu'au \omterm{def:s1:index-rebalancing:fund}{jour de classement} et atteint son maximum pour les actions proches du seuil de coupure.
\end{solution}

\begin{solution}{iq:s1:index-rebalancing:4}
Les prévisions erronées~; une position encombrée qui a déjà intégré le changement~; des \omterm{def:s1:index-rebalancing:fund}{fonds indiciels} qui échangent avant ou après la clôture~; une baisse du marché avant le \omterm{def:s1:index-rebalancing:fund}{jour de classement}
qui déplace le seuil de coupure~; des actions proches du seuil qui sont petites et peu liquides.
\end{solution}

\begin{solution}{iq:s1:index-rebalancing:5}
Si le déséquilibre provient d'une demande indicielle mécanique et que le cours reflète déjà un fort positionnement préalable, il faut vendre dans ce mouvement et prévoir de racheter à mesure que la pression s'inverse~;
dimensionner la position selon la profondeur de l'enchère et selon la part du mouvement déjà réalisée avant la clôture.
\end{solution}

\begin{solution}{iq:s1:index-rebalancing:6}
Si les arbitragistes achètent $A$ avant l'annonce et que les \omterm{def:s1:index-rebalancing:fund}{fonds indiciels} achètent $D$ à la clôture effective, le marché absorbe $A$ tôt et $D - A$ plus tard, et revend $A$
aux fonds~; avec des poussées proportionnelles à la demande nette, le mouvement précoce est proportionnel à $A$ et le mouvement postérieur à l'annonce à $D - A$, de sorte que le rendement
après l'annonce est proportionnel à $1 - A/D$ et s'annule à mesure que le capital d'arbitrage s'approche de la demande indicielle.
\end{solution}
